\chapter{Lời hứa có biên và các điều kiện bảo đảm}\label{ch:phuongphap}

Chương này trả lời RQ1 bằng cách xác định đại lượng được ghi sổ, thời điểm kiểm tra và phạm vi
của chứng nhận. Từ đó, chương xây dựng khung bốn họ cổng để trả lời phần lý thuyết của RQ3.
Các bảo đảm đều có giả thiết; chúng không thay thế việc đo chất lượng dịch vụ và các biến động
mà hành khách thực sự nhìn thấy ở Chương~\ref{ch:thucnghiem}.

\section{Tổng quan cơ chế}

RideBound công bố giờ đón và giờ trả dự kiến cho khách đã nhận. Mỗi khi công bố lại, độ dời
được ghi vào sổ cái riêng của khách. Sổ cái không giảm; phương án làm vượt trần cứng bị loại
trước khi bộ chọn đội xe giải bài toán quy hoạch nguyên.

Điểm cần lựa chọn là độ dời nào được tính. Khóa luận so phương án mới với tuyến giữ nguyên
trên cùng trạng thái hiện tại. Phần chênh giữa hai phép chiếu được quy cho quyết định. Phần
chênh giữa tuyến giữ nguyên và giờ đã báo là phần thay đổi khi giữ tuyến. Một trần trên phần
thứ nhất không đồng nghĩa trần trên toàn bộ biến động khách nhìn thấy.

Một quyết định có miễn trừ vẫn có thể tiếp tục phục vụ khách nhưng phải ghi vi phạm và không
được gắn nhãn vận hành bình thường. Đây là một phần của phạm vi chứng nhận, không phải một
cách biến lời hứa đã vượt trần thành lời hứa được giữ.

\section{Mô hình bài toán mở rộng}

\subsection{Lời hứa công bố}

Mở rộng mô hình bài toán ở Chương 2, ta bổ sung khái niệm lời hứa. Với mỗi hành khách $r$ đã được xác nhận,
hệ thống duy trì một \emph{lời hứa đang mở}, gồm định danh xe được phân công, điểm đón, điểm
trả, giờ đón dự kiến và giờ trả dự kiến. Lời hứa này được công bố cho hành khách và được ghi
vào biên bản. Mỗi khi hệ thống ra quyết định, lời hứa có thể được công bố lại với giá trị mới.
Trong chế độ xác nhận đặt xe dùng cho FleetPy, đề nghị tạm thời chưa mở lời hứa; mốc đầu
được ghi khi khách xác nhận đặt xe.

Chính sách phân biệt ba pha của yêu cầu: \emph{đã nhận} (hệ thống đã chấp nhận yêu cầu),
\emph{đang chờ đón}, và \emph{đang trên xe}. Việc phân biệt ba pha
là cần thiết vì các chiều cam kết có ý nghĩa khác nhau ở từng giai đoạn --- chẳng hạn giờ đón
không còn ý nghĩa sau khi khách đã lên xe.

\subsection{Các chiều cam kết}

Hệ thống theo dõi mười chiều cam kết, chia thành ba nhóm. Nhóm thời gian gồm tổng độ dời giờ
đón và tổng độ dời giờ trả, tính bằng mili giây. Nhóm định danh gồm số lần đổi xe, số lần đổi
điểm đón, số lần đổi điểm trả. Nhóm cấu trúc gồm khoảng cách dời điểm đón và điểm trả tính bằng
mili-mét, số cặp điểm dừng của khách cũ bị đảo thứ tự, số điểm dừng chèn thêm trước khi đón, và số lần sửa
đổi đáng kể theo ngưỡng được cấu hình.

Mỗi chiều có thể được đặt một trần cứng, hoặc để không giới hạn. Trong thực nghiệm, cấu hình
được dùng nhiều nhất đặt trần cho chiều \emph{tổng độ dời giờ trả} và khóa cứng ba chiều định
danh về không, tức cấm hoàn toàn việc đổi xe hay đổi điểm dừng sau khi đã hứa.

\section{Phép phân rã ba mốc}
\label{muc:bamoc}

\subsection{Định nghĩa}
\label{muc:dinhnghia-bamoc}

Tại một thời điểm quyết định, với một tọa độ thời gian của hành khách $r$ đang có lời hứa mở,
hệ thống dựng ba mốc:

\begin{itemize}[leftmargin=1.2cm,itemsep=2pt,topsep=3pt]
\item $v$ --- \textbf{previous}: lời hứa đã công bố lần gần nhất cho hành khách;
\item $e$ --- \textbf{exogenous}: chính kế hoạch đó, chiếu lại từ vị trí thật hiện tại của xe.
Cùng lộ trình, cùng phân công, chỉ khác điểm xuất phát và thời điểm chiếu;
\item $p$ --- \textbf{proposed}: lời hứa mà phương án đang xét sẽ tạo ra.
\end{itemize}

Từ ba mốc này định nghĩa ba đại lượng:

\begin{align}
\text{độ dời do quyết định} \quad & A = |p - e|\\
\text{thay đổi khi giữ tuyến} \quad & d = |e - v|\\
\text{độ dời khách nhìn thấy} \quad & B = |p - v|
\end{align}

Mốc $e$ trả lời câu hỏi trong mô hình: nếu nhà điều hành giữ tuyến, giờ hẹn hiện được chiếu
là bao nhiêu? Chênh lệch với giờ đã hứa là phần thay đổi khi giữ tuyến, gồm cả giao thông,
tiến độ thực và độ trễ thực thi (mục~\ref{muc:phanra}). Chênh lệch với phương án đang xét
là độ dời quy cho quyết định trong phép đối chiếu tại cùng trạng thái.

\paragraph{$A$ là đại lượng đối xứng.} Cài đặt tính $A = |p - e|$ theo đúng trị tuyệt đối
(\texttt{PromiseDeltaCalculator}, hàm \texttt{AbsoluteDifference}): một quyết định kéo lời hứa
\emph{sớm} hơn mốc giữ tuyến cũng bị tính phí bằng đúng độ lệch đó, giống hệt một quyết định
đẩy muộn. Không có quy ước một chiều nào trong mã.

Cơ sở hành vi vì vậy có thể tính phí cả một sửa chữa khôi phục giờ đã hứa. Xét
$v = 600$, $e = 720$, $p = 600$ (giây): giao thông đẩy lời hứa muộn đi 120 giây, và phương án
đang xét kéo nó về đúng giờ đã hứa. Khi đó $A = 120$ trong khi $B = 0$: khách không thấy giờ
dự kiến thay đổi, nhưng sổ cái hành vi vẫn ghi thêm 120 giây. Một trần hành vi chặt có thể
loại phương án khôi phục đúng giờ đã hứa.

Đây là giới hạn của sổ cái đối xứng. Bộ sinh chèn thuần trong các phép so chính chưa dùng
sửa chữa kéo giờ hẹn về sớm hơn mốc giữ tuyến. Kết quả ấy là đặc điểm của bộ sinh và dữ liệu,
không làm mất phản ví dụ. Một sổ cái chỉ tính phần gây hại cần định nghĩa và đánh giá riêng.

\subsection{Ví dụ minh họa}

Một hành khách đặt xe lúc 9:00 và được báo sẽ được đón lúc 9:10. Đến 9:05, hệ thống cân nhắc
thay đổi kế hoạch vì có khách mới. Giả sử tại thời điểm đó xe đang chậm hơn dự kiến 40 giây.

\begin{figure}[htbp]
\centering
\includegraphics[width=0.95\linewidth]{hinh/co-che-ba-moc.pdf}
\caption{Ví dụ dựng để giải thích ba mốc: tuyến giữ nguyên trễ 40 giây, phương án mới thêm
30 giây. Sơ đồ này không phải một quan sát thực nghiệm.}
\label{hinh:bamoc}
\end{figure}
Hình~\ref{hinh:bamoc} dùng giờ đã báo làm gốc để minh họa phép đối chiếu.

Với phương án đổi kế hoạch, cơ sở hành vi tính $9{:}11{:}10 - 9{:}10{:}40 = 30$ giây, còn cơ sở
kết quả tính $9{:}11{:}10 - 9{:}10{:}00 = 70$ giây. Khác biệt 40 giây chính là phần xe chạy chậm.

Bây giờ xét phương án \emph{không làm gì}, tức $p = e = 9{:}10{:}40$. Cơ sở hành vi tính 0 giây,
nên riêng cổng hành vi không loại phương án này nếu sổ cái trước đó hợp lệ. Cơ sở kết quả tính 40 giây --- đã vượt trần 30 giây. Theo cách
tính thứ hai, \emph{ngay cả việc không làm gì cũng là vi phạm}. Nếu mọi phương án khác cũng vượt,
chiếc xe không còn kế hoạch hợp lệ nào và hệ thống dừng. Đây là hạt nhân của toàn bộ phần lý
thuyết ở mục \ref{muc:haicosotinh}.

\section{Sổ cái cam kết}

Với mỗi cặp (hành khách, chiều cam kết), hệ thống duy trì một bộ đếm cộng dồn. Mỗi lần công bố
lại lời hứa, bộ đếm được cập nhật

\begin{equation}
L_{r,\text{sau}} = L_{r,\text{trước}} + \delta_r
\end{equation}

trong đó $\delta_r$ là độ dời của lần công bố ấy theo cơ sở đang dùng. Vì $\delta_r \ge 0$, bộ
đếm đơn điệu tăng.

Một phương án được coi là hợp lệ về mặt cam kết nếu với mọi hành khách $r$ bị ảnh hưởng và mọi
chiều $k$ có đặt trần, ta có $L_{r,k,\text{sau}} \le \beta_k$. Phương án vi phạm bị loại khỏi
tập lựa chọn trước khi bài toán quy hoạch nguyên được giải, chứ không bị phạt trong hàm mục
tiêu.

Ba đặc điểm sau xác định chính xác cơ chế được kiểm trong khóa luận. Thứ
nhất, nó \emph{theo từng hành khách}: không có phép cộng gộp nào làm loãng người chịu thiệt
nặng nhất. Thứ hai, nó \emph{cộng dồn theo thời gian}: một hành khách bị dời năm lần, mỗi lần
mười giây, tiêu tốn đúng bằng một hành khách bị dời một lần năm mươi giây. Thứ ba, nó là
\emph{ràng buộc cứng}: phương án vượt trần bị loại trong vận hành bình thường.

\section{Hai cơ sở tính, các định lý và khung bốn họ cổng}
\label{muc:haicosotinh}

\subsection{Phát biểu hai cơ sở}

Khóa luận xét hai cơ sở tính sổ cái từ ba mốc ở mục~\ref{muc:bamoc}. Hai cơ sở này không
tương đương và đều tính cả độ dời sớm lẫn độ dời muộn.

\begin{dinhnghia}[Cơ sở hành vi]
Sổ cái cộng dồn đại lượng $A = |p-e|$, tức độ dời so với phép chiếu khi giữ tuyến.
Trong cài đặt, đây là \texttt{CommitmentBudgetBasis.DecisionInduced}. 
\end{dinhnghia}

\begin{dinhnghia}[Cơ sở kết quả]
Sổ cái cộng dồn đại lượng $B = |p-v|$, tức độ dời so với lời hứa đã công bố gần nhất. Trong cài
đặt, đây là \texttt{CommitmentBudgetBasis.CustomerVisible}. 
\end{dinhnghia}

\subsection{Phương án giữ tuyến trong danh mục ứng viên}

Danh mục ứng viên của mỗi xe phải chứa đúng một phương án giữ nguyên kế hoạch hiện tại,
còn gọi là phương án không làm gì (no-op). Bộ sinh, bộ giữ danh mục và bộ đánh giá kiểm yêu
cầu này. Bài toán chọn buộc mỗi xe chọn \emph{đúng} một phương án, kể cả khi không nhận thêm
khách. Việc giữ một ứng viên no-op trong danh mục chưa bảo đảm nó qua mọi cổng khả thi.

Với phương án không làm gì, kế hoạch không đổi, nên phép chiếu của nó từ vị trí hiện tại
\emph{chính là} mốc $e$. Tức là $p(\text{no-op}) = e$.

\subsection{Bất đẳng thức ba mốc và điều kiện dấu bằng}
\label{muc:phanra}

Ba đại lượng liên hệ với nhau qua bất đẳng thức tam giác:

\begin{equation}
\boxed{\,B \le A + d\,}
\label{eq:phanra}
\end{equation}

Dấu bằng xảy ra khi và chỉ khi hai hiệu $(p-e)$ và $(e-v)$ không trái dấu. Khóa luận không
giả định điều kiện đó luôn đúng; độ hụt so với dấu bằng được báo ở mục~\ref{muc:quytrach}.

\paragraph{Vì sao dấu bằng có thể hỏng.} Cả ba đại lượng đều là trị tuyệt đối. Nếu phần thay đổi
khi giữ tuyến kéo $e$ \emph{sớm} hơn $v$ trong khi quyết định đẩy $p$ \emph{muộn} hơn $e$, hai hiệu
ngược dấu triệt tiêu một phần và $|p-v| < |p-e| + |e-v|$. Vì thế tổng phần do quyết định và
tổng phần giữ tuyến không nhất thiết bằng tổng biến động nhìn thấy, dù cùng đo bằng giây.

\paragraph{Phạm vi của phép đối chiếu.}
Mốc $e$ được dựng trên trạng thái sau khi lô sự kiện đã được áp dụng. Vì vậy $d$ có thể gồm
cập nhật thời gian đi, tiến độ thực, thời gian dừng và độ trễ thực thi. Không được gọi nó là
tác động thuần của giao thông. Tuyến đã chọn trước đó cũng ảnh hưởng vị trí xe hiện tại,
nên $d$ không được giả định độc lập với chính sách.

Phép đối chiếu là một cách hạch toán trong mô hình, tại cùng trạng thái, không phải một ước
lượng tác động nhân quả tổng thể của việc đổi chính sách. Dấu bằng trong~\eqref{eq:phanra}
phải được kiểm, không được giả định. Các tỉ số giữa tổng phần do quyết định và tổng biến
động nhìn thấy ở mục~\ref{muc:quytrach} có tử số, mẫu số rõ ràng; chúng không phải các tỉ
trọng của một phân rã cộng tính luôn đúng.

\paragraph{Điều kiện ngân sách trước quyết định (B).}
Sổ cái trước quyết định không vượt trần đang áp dụng. Điều kiện này được duy trì khi sổ cái
khởi tạo bằng không, trần cố định suốt chuỗi và các lần công bố trước đều qua cổng trong
vận hành bình thường. Các kết quả giữ tuyến dưới đây cần điều kiện (B); chúng không áp
dụng nguyên trạng cho sổ cái đã vượt trần sau miễn trừ hoặc vừa đổi pha áp dụng trần.

\paragraph{Điều kiện các kiểm tra khác (N).}
Phương án giữ tuyến tại kỷ nguyên đang xét qua mọi kiểm tra ngoài cổng thời gian đang so,
gồm cấu trúc tuyến, sức chứa, lịch trình, khóa và vòng đời. Cổng ngân sách không thay thế
tính khả thi vật lý. Hai điều kiện (B) và (N) phải được kiểm riêng.

\subsection{Ba định lý}


Các định lý xét một tọa độ thời gian, cùng trần cố định và cổng không có miễn trừ. Khi có
nhiều cổng, các điều kiện tương ứng phải đúng cho từng cổng đang áp dụng.

\begin{dinhly}[Cổng hành vi giữ tuyến dưới điều kiện ngân sách]

Dưới điều kiện (B), cổng ngân sách hành vi không loại phương án giữ tuyến. Nếu thêm (N),
tập phương án sau các kiểm tra không rỗng.
\end{dinhly}

\begin{proof}
Với phương án không làm gì, $A = |p - e| = |e - e| = 0$. Do đó sổ cái sau bằng sổ cái trước,
mà theo (B), sổ cái trước không vượt $\beta$. Vậy cổng ngân sách không loại phương án giữ
tuyến; theo (N), phương án đó cũng qua các kiểm tra khác và là một phần tử của tập khả thi.
\end{proof}

Kết quả dựa vào đẳng thức $p=e$ của phép chiếu khi giữ tuyến, không phải khả năng kiểm soát
giao thông. Hệ thống vẫn có thể dừng vì một kiểm tra khác hoặc vì (B) không còn đúng.

\begin{dinhly}[Cơ sở kết quả có thể vô nghiệm]

Dưới cơ sở kết quả, tồn tại trạng thái thỏa (B) và (N) mà cổng ngân sách làm rỗng tập
phương án của một xe.
\end{dinhly}

\begin{proof}
Xét một xe chở một khách cũ, không có yêu cầu mới và chỉ có phương án giữ tuyến. Cho sổ cái
trước bằng $0$, trần $\beta=30$ giây và phép chiếu giữ tuyến muộn hơn lời hứa trước $40$ giây.
Chọn các cửa sổ vật lý đủ rộng để phương án vẫn thỏa (N). Khi đó (B) đúng nhưng
$B=|e-v|=40>\beta$, nên phương án duy nhất bị loại. Đây là một trạng thái dựng, không phải
một quan sát thực nghiệm.
\end{proof}

\begin{dinhly}[Bao hàm quỹ đạo dưới điều kiện (H)]
\label{dl:tapcon}
Xét cùng kịch bản và chân trời hữu hạn, cùng phép chiếu, quy tắc chuyển trạng thái và các
ràng buộc khác; hai sổ cái đều khởi tạo bằng không. Giả sử điều kiện \emph{(H)} đúng ở mọi
lần công bố của mọi quỹ đạo khả thi dưới cơ sở kết quả:
$|p-e|\le|p-v|$. Khi đó, với cùng trần $\beta$, mọi quỹ đạo ấy cũng khả thi dưới cơ sở
hành vi. Vì vậy số khách phục vụ lớn nhất trên toàn bộ quỹ đạo của cơ sở kết quả không
vượt số khách phục vụ lớn nhất của cơ sở hành vi.
\end{dinhly}

\begin{proof}
Giữ nguyên một quỹ đạo khả thi dưới cơ sở kết quả và ghi hai sổ cái trên chính quỹ đạo đó.
Ký hiệu $L^{A}_k$, $L^{B}_k$ là các giá trị sau $k$ lần công bố, với
$L^{A}_0=L^{B}_0=0$. Nếu $L^{A}_k\le L^{B}_k$, thì theo (H) và cổng kết quả,
$$L^{A}_k + A \;\le\; L^{B}_k + B \;\le\; \beta,$$
nên phương án đó cũng được cơ sở hành vi chấp nhận, và bất đẳng thức giữa hai sổ cái được bảo
toàn sang bước $k+1$. Vậy mọi chuỗi quyết định hợp lệ dưới cơ sở kết quả cũng hợp lệ dưới cơ sở
hành vi, tức tập quỹ đạo khả thi của cơ sở kết quả là \emph{tập con} của tập quỹ đạo khả thi của
cơ sở hành vi. Lấy cực đại số khách phục vụ trên hai tập lồng nhau cho kết luận cuối.
\end{proof}

\paragraph{Phạm vi của quan hệ bao hàm.}
Điều kiện (H) không là hệ quả của bất đẳng thức tam giác. Khi phần thay đổi giữ tuyến hấp thụ
một phần độ dời do quyết định, $A>B$ có thể xảy ra. Chẳng hạn ví dụ $v=600$, $e=720$,
$p=600$ ở mục~\ref{muc:dinhnghia-bamoc} cho $A=120>B=0$. Việc dấu bằng của bất đẳng thức
tam giác không đạt chưa đủ để kết luận (H) hỏng, và kiểm trên phương án đã chọn chưa chứng
minh (H) cho toàn bộ tập quỹ đạo. Vì vậy định lý không được dùng như kết luận về panel.

Định lý so tập quỹ đạo và chính sách tối ưu toàn quỹ đạo. Bộ chọn đang cài chỉ tối ưu ở từng
kỷ nguyên, không nhìn trước yêu cầu tương lai. Chọn một phương án khác hôm nay có thể làm
trạng thái ngày mai khác đi, nên một cấu hình chặt hơn vẫn có thể tình cờ hoàn tất nhiều
khách hơn ở một ô. Phần thực nghiệm phải kiểm riêng, không suy trực tiếp từ định lý.

\begin{hequa}[Dải khác biệt tại một lần công bố]

Xét một lần công bố mà ngay trước đó hai sổ cái bằng nhau, $L^{A} = L^{B} = L$. Một phương án được
cơ sở hành vi giữ lại nhưng bị cơ sở kết quả loại khi và chỉ khi $L + A \le \beta < L + B$, tức
$L + A$ rơi vào dải $(\beta - (B - A),\, \beta]$. Vì $B \le A + d$, bề rộng dải không quá $d$ của
lần công bố đó, và đúng bằng $d$ khi dấu bằng trong \eqref{eq:phanra} xảy ra; khi (H) hỏng
($B < A$) dải rỗng.
\end{hequa}

Hệ quả chỉ so một lần công bố khi hai sổ cái bằng nhau. Trên cả chuỗi, sổ cái có thể tách
nhau và các chính sách đi quỹ đạo khác nhau. Nó không cho một công thức dự báo chênh số khách.

\subsection{Khung so sánh bốn họ cổng}


Ba định lý trên so hai cơ sở tính của \emph{cùng một} cơ chế. Các công trình gần nhất lại dùng
những ràng buộc lời hứa có hình thức khác: Schulz và Pfeiffer \cite{schulz2026} đặt một hạn chót
tính từ lời hứa đầu tiên; một cửa sổ quanh lời hứa đầu gần với ``hộp tĩnh'' gồm cửa sổ thời gian
và thời gian đi tối đa của Sharif Azadeh và cộng sự \cite{sharifazadeh2022}. Mục này đặt tất cả
vào cùng một khung để so được, và rút ra kết quả trung tâm của chương.

\paragraph{Đại lượng có dấu và đại lượng cộng dồn.} Xét một hành khách và tọa độ giờ trả. Tại mỗi
lần công bố $t$ đặt
\begin{equation}
x_t = e_t - v_t, \qquad c_t = p_t - e_t, \qquad z_t = p_t - v_t,
\end{equation}
khi đó $d_t = |x_t|$, $A_t = |c_t|$, $B_t = |z_t|$, và đồng nhất \emph{có dấu} $z_t = x_t + c_t$
luôn đúng --- bất đẳng thức \eqref{eq:phanra} chỉ xuất hiện khi lấy trị tuyệt đối. Gọi $p_0$ là
lời hứa \emph{đầu tiên} của hành khách. Cộng dồn trên cùng lưới công bố:
\begin{gather}
C = \sum_t A_t, \qquad X = \sum_t d_t, \qquad V = \sum_t B_t, \qquad E = |p_T - p_0|,\\
M = \max_t |p_t - p_0|, \qquad M^{+} = \max_t \max(0,\, p_t - p_0).
\end{gather}
$E$ là độ lệch hai đầu, $M^{+}$ là phần \emph{muộn} cực đại so với lời hứa đầu --- đúng đại lượng mà
một hạn chót kiểu Schulz--Pfeiffer chặn. Trong mục này $\kappa$ ký hiệu trần chung của một họ ràng
buộc; với hai sổ cái của khóa luận, $\kappa$ chính là trần $\beta$.

\begin{menhde}[Chuỗi đại lượng]
\label{md:chuoi}
Nếu chuỗi công bố liên tục, tức mỗi lần công bố so với lần ngay trước nó ($v_t = p_{t-1}$), thì
\[
\max(0,\, p_T - p_0) \;\le\; M^{+} \;\le\; M \;\le\; V \;\le\; C + X
\qquad\text{và}\qquad E \le M .
\]
\end{menhde}

\begin{proof}
Bất đẳng thức đầu và $E \le M$ đúng vì giá trị tại lần công bố cuối không vượt cực đại trên mọi
lần công bố; $M^{+} \le M$ vì $\max(0,\cdot) \le |\cdot|$. Với chuỗi liên tục,
$p_t - p_0 = \sum_{s \le t} z_s$, nên $|p_t - p_0| \le \sum_{s \le t} |z_s| \le V$. Cuối cùng
$|z_t| \le |c_t| + |x_t|$ ở từng lần công bố, cộng dồn lại được $V \le C + X$.
\end{proof}

Bản thân chuỗi bất đẳng thức này là kiến thức cơ bản, \emph{không} phải đóng góp; nó chỉ là công
cụ cho định lý dưới đây.

\paragraph{Bốn họ cổng.} Bảng \ref{bang:hocong} xếp bốn cách ràng buộc lời hứa theo đại lượng mà
chúng kiểm ở mỗi quyết định. (Schulz và Pfeiffer không công bố lại lời hứa --- khách chỉ được báo
giờ tới và hạn chót một lần --- nhưng hạn chót được kiểm ở mọi lần chèn khách mới về sau.)

\begin{table}[htbp]
\centering
\caption{Bốn họ cổng cam kết với trần $\kappa$}
\label{bang:hocong}
\small
\begin{tabularx}{\linewidth}{@{}llXX@{}}
\toprule
\textbf{Họ} & \textbf{Kiểm ở mỗi quyết định} & \textbf{Tên gọi} & \textbf{Ví dụ}\\
\midrule
$C$ & $\sum_t A_t \le \kappa$ & theo hành vi & cơ sở hành vi của khóa luận\\
$M^{+}$ & $p_t - p_0 \le \kappa$ & hạn chót một chiều & Schulz và Pfeiffer \cite{schulz2026}\\
$M$ & $|p_t - p_0| \le \kappa$ & hạn chót hai chiều & cửa sổ quanh lời hứa đầu\\
$V$ & $\sum_t B_t \le \kappa$ & theo kết quả & cơ sở kết quả của khóa luận\\
\bottomrule
\end{tabularx}
\end{table}

\paragraph{Phạm vi khả thi.}
Điều kiện (N) vẫn cần cho khung bốn họ cổng. Nếu xe tới đón ngoài cửa sổ đã chấp nhận,
một bất biến vòng đời có thể làm lần chạy dừng. Mục~\ref{muc:trongdiem} phân loại các lần
dừng theo nguyên nhân.

Việc danh mục phương án yêu cầu tuyến giữ nguyên cũng là một lựa chọn cài đặt. Trong chiều
trôi thuận lợi, một cổng có thể loại tuyến giữ nguyên nhưng vẫn giữ phương án khác; không
được coi mọi lần tuyến giữ nguyên bị loại là một chứng minh tập toán học hoàn toàn rỗng.

\paragraph{Quỹ đạo dự phòng.} Xét chuỗi hữu hạn từ lời hứa đầu tiên, khởi tạo sổ cái bằng
không, mà ở mọi kỷ nguyên xe giữ nguyên tuyến. Trên đó
$p_t = e_t$, nên $c_t = 0$ với mọi $t$, và
\[
C_{\text{noop}} = 0, \qquad z_t = x_t, \qquad V_{\text{noop}} = X_{\text{noop}} ,
\]
trong đó chỉ số ``noop'' chỉ đại lượng tính dọc quỹ đạo dự phòng (khác với quỹ đạo thật đã chạy).

\begin{dinhly}[Thứ tự ngưỡng giữ tuyến của các họ cổng]
\label{dl:hocong}
Xét một hành khách, một tọa độ thời gian và trần $\kappa$ cố định. Giả sử (N) tại các
kỷ nguyên đang xét.
\begin{enumerate}[label=(\alph*),leftmargin=1cm,itemsep=2pt,topsep=3pt]
\item \emph{Tại một quyết định.} Nếu (B) đúng, cổng $C$ giữ phương án không làm gì với mọi
$\kappa\ge0$. Vì vậy cổng này không làm rỗng tập phương án tại quyết định đó.
\item \emph{Quỹ đạo dự phòng.} Trên quỹ đạo dự phòng với chuỗi công bố liên tục, ngưỡng nhỏ nhất
bảo đảm phương án không làm gì không bị loại ở mọi kỷ nguyên là
\[
\kappa^{*}_{C} = 0 \;\le\; \kappa^{*}_{M^{+}} = M^{+}_{\text{noop}} \;\le\;
\kappa^{*}_{M} = M_{\text{noop}} \;\le\; \kappa^{*}_{V} = X_{\text{noop}} ,
\]
tức ba họ còn lại có thể loại phương án không làm gì, với ngưỡng giữ tuyến phụ thuộc phần
trôi dạt. $M^{+}$ chỉ xét độ muộn so với lời hứa đầu (Hệ quả~\ref{hq:chieutroi}).
\end{enumerate}
\end{dinhly}

\begin{proof}
(a) Sổ cái trước kỷ nguyên $t$ thỏa $C_{t-1} \le \kappa$ vì lần công bố trước đã qua cổng. Phương
án không làm gì có $c_t = 0$ nên $C_t = C_{t-1} \le \kappa$: cổng không loại nó, và theo (N) nó
cũng qua mọi kiểm tra khác, nên tập phương án không bị làm rỗng \emph{vì cổng}. (b) Phương án không
làm gì bị loại ngay tại lần công bố đầu tiên mà đại lượng được kiểm vượt $\kappa$, nên ngưỡng nhỏ
nhất đúng bằng giá trị lớn nhất của đại lượng đó dọc quỹ đạo dự phòng (với $V$ là giá trị cuối, vì
$V$ cộng dồn đơn điệu). Thứ tự giữa các ngưỡng là Mệnh đề \ref{md:chuoi} áp cho quỹ đạo có
$C = 0$; bước $M \le V$ cần chuỗi công bố liên tục.
\end{proof}

Phần (a) dựa vào $p=e$ của tuyến giữ nguyên và ngân sách trước đó hợp lệ. Phần (b) so các
ngưỡng trên cùng quỹ đạo dự phòng, không so số khách của các nhánh đã chọn tuyến khác nhau.
Chứng minh sơ cấp; khung này có ích vì phân biệt các lời hứa cùng đơn vị giây và tạo dự
đoán kiểm được từ chuỗi công bố. Trong bốn họ đang xét, chỉ $C$ không tính phần trôi vào
ngân sách. Điều đó không phải bảo đảm hệ thống luôn chạy trọn hay khách không thấy đổi giờ.

\begin{hequa}[Khi thời gian đi tất định]

Trên quỹ đạo dự phòng thỏa (N), nếu phép chiếu và tiến độ thực thi nhất quán để
$x_t\equiv0$, và chuỗi công bố liên tục, thì
$M^{+}_{\text{noop}} = M_{\text{noop}} = X_{\text{noop}} = 0$ và cả bốn họ đều không bao giờ loại
phương án không làm gì, với mọi $\kappa \ge 0$.
\end{hequa}

Thời gian đi cố định chưa đủ để suy ra $x_t=0$: tiến độ thực thi và phép chiếu cũng phải
nhất quán. Panel mạng tĩnh của khóa luận vẫn có trôi nhỏ do mô phỏng rời rạc. Vì vậy hệ quả
là một trường hợp lý tưởng của mô hình; việc áp dụng cần kiểm biên bản.

\begin{hequa}[Phụ thuộc chiều trôi]
\label{hq:chieutroi}
Trên quỹ đạo dự phòng với chuỗi công bố liên tục, $p_t - p_0 = \sum_{s \le t} x_s$, nên
\[
M^{+}_{\text{noop}} = \max_t \Big(\sum_{s \le t} x_s\Big)^{+}, \qquad
M_{\text{noop}} = \max_t \Big|\sum_{s \le t} x_s\Big|, \qquad
V_{\text{noop}} = \sum_t |x_t| .
\]
Nếu trôi \emph{bất lợi} (mọi $x_s\ge0$, phép chiếu giờ đến dần muộn hơn) thì
$M^{+}_{\text{noop}} = M_{\text{noop}} = V_{\text{noop}}$: ba họ $M^{+}$, $M$, $V$ cùng loại phương
án không làm gì ở cùng một ngưỡng. Nếu trôi \emph{thuận lợi} (mọi $x_s\le0$, phép chiếu
giờ đến dần sớm hơn) thì
$M^{+}_{\text{noop}} = 0$: hạn chót một chiều không bao giờ loại phương án không làm gì, trong khi
$M_{\text{noop}} = V_{\text{noop}}$ vẫn có thể lớn, nên hạn chót hai chiều và cổng kết quả vẫn có
thể loại nó. Cổng hành vi không loại nó ở cả hai chiều.
\end{hequa}

Loại phương án không làm gì chưa đủ để kết luận tập phương án rỗng. Trên quỹ đạo dự phòng
chỉ trôi bất lợi, nếu mọi phương án chèn có $c\ge0$, thì chúng đều muộn ít nhất bằng
phương án giữ tuyến tại cùng trạng thái. Khi một trong các cổng $M^{+}$, $M$, $V$ loại
phương án giữ tuyến, nó cũng loại các phương án chèn này. Với trôi thuận lợi, phương án
chèn làm muộn lại có thể lọt vào cửa sổ $M$ hay dưới trần $V$. Khi đó tập toán học vẫn có
thể còn phần tử, dù cài đặt đòi ứng viên no-op và dừng khi ứng viên này mất tính hợp lệ.

Trên hệ thật, phần giữ tuyến có thể đổi dấu và không chỉ gồm giao thông. Dự đoán theo ô
phải tính trên chuỗi đo được của từng khách. Khi trôi thuận lợi, hạn chót một chiều có thể
cho phép chèn thêm bằng phần chùng vừa tạo ra, còn cổng hành vi vẫn tính đủ độ dời do quyết
định. Đây là khả năng đánh đổi, không phải ưu thế dịch vụ của một họ đã được chứng minh.

\subsection{Khi nào chặn theo chuỗi khác chặn hai đầu}


\begin{menhde}[Bộ sinh chèn thuần]

Xét một chuỗi công bố liên tục. Nếu mọi phương án được chọn đều có $p \ge e$ với mọi hành khách
cũ, tức $c_t \ge 0$ ở mọi lần công bố, thì $C = \sum_t c_t$. Nếu thêm $x_t \equiv 0$ thì
$C = V = p_T - p_0 = M^{+} = M$: trần theo chuỗi, trần theo hai đầu và trần theo hành vi là một.
Khi $x_t \not\equiv 0$, kết luận thứ hai không còn được bảo đảm.
\end{menhde}

\begin{proof}
$c_t \ge 0$ cho $A_t = c_t$. Khi $x_t \equiv 0$ thì $z_t = c_t \ge 0$; với chuỗi liên tục
$p_t - p_0 = \sum_{s \le t} z_s$, nên $p_t$ đơn điệu không giảm, $p_T - p_0 = \sum_t z_t = V = C$,
và cực đại của $p_t - p_0$ đạt ở lần công bố cuối.
\end{proof}

Kết quả này giới hạn cách phát biểu đóng góp. Khi $c_t\ge0$, ta có
$|C-(p_T-p_0)|=|\sum_t x_t|\le X$, nên trần cộng dồn hành vi gần trần độ muộn hai đầu
khi $X$ nhỏ trên cùng chuỗi. Khóa luận không coi hai hình thức khác nhau là một lợi thế
đã được chứng minh độc lập với bộ sinh.

Hai ví dụ dựng tách hiệu hai đầu khỏi độ dài chuỗi. Chuỗi $600\to660\to600$ có $E=0$ nhưng
$V=120$ giây; chuỗi $600\to660\to720$ có $E=V=120$ giây. Ghi nhận dao động không đồng nghĩa
đã đo được tổn thất hài lòng; khóa luận không có dữ liệu lựa chọn hay cảm nhận của hành khách.

\section{Các khóa và cơ chế bảo vệ bổ sung}

Khóa cứng cấm đổi một trường lời hứa trong pha được cấu hình. Hàng rào đóng băng cấm đổi khi
giờ đón còn cách hiện tại dưới một ngưỡng. Cơ chế một chiều cho phép giờ dự báo tốt lên nhưng
cấm xấu đi so với mốc tham chiếu được chọn. Đây là những hình thức có tiền lệ trong văn liệu,
không phải các khái niệm mới của khóa luận.

Các cấu hình này phải được đọc cùng bộ sinh phương án. Nếu bộ sinh không thể làm giờ hẹn
tốt lên, quyền cho phép cải thiện không tạo thêm lựa chọn. Ngược lại, một bộ sinh biết sửa
tuyến có thể làm cơ chế một chiều khác với khóa hai chiều. Mục~\ref{muc:kiemtraphu} báo các
quan sát phụ trong phạm vi cài đặt đã thử.

\section{Chứng nhận và phục hồi có miễn trừ}
\label{muc:phuchoi}

\subsection{Phạm vi của chứng nhận}

``Được chứng nhận'' ở đây nghĩa là quyết định có bằng chứng kiểm tra gắn với phiên bản,
cấu hình và lịch sử lời hứa. Chứng cứ ghi giá trị trước, độ dời, giá trị sau và các lý do
loại hoặc miễn trừ. Nó không phải chứng nhận pháp lý, bảo đảm thời gian đến thật, hay chứng
minh hình thức tính đúng của toàn bộ chương trình.

Quyết định có miễn trừ được đánh dấu \texttt{normalOperation=false}. Nó có thể được thực thi
để tiếp tục phục vụ khách đã nhận, nhưng không được tính là giữ lời hứa của cổng bị vượt.
Tỉ lệ quyết định có miễn trừ và tỉ lệ khách vượt lời hứa có mẫu số khác nhau, phải báo riêng.

\subsection{Hai quy tắc phục hồi}

\textbf{Quy tắc (a), giữ tuyến.} Xe có tuyến tham chiếu đã vượt ngân sách hoặc hạn chót vẫn
đi tuyến đó, với vi phạm được ghi nhận có kiểu. Quy tắc có thể làm giảm cơ hội nhận khách
mới vì xe bị hạn chế đổi tuyến. Việc tiếp tục chạy không biến vi phạm thành tuân thủ.

\textbf{Quy tắc (b), không tệ hơn tham chiếu.} Xe bị vướng được đổi tuyến trên cùng xe chỉ
khi từng cổng được xin miễn trừ cũng bị tuyến giữ nguyên vi phạm và phương án mới không
tệ hơn ở đúng cổng ấy. Với cận muộn, giờ trả mới không muộn hơn mốc giữ tuyến; với cận sớm,
nó không sớm hơn. Với chiều ngân sách, giá trị sau không lớn hơn giá trị sau của tham chiếu
đã vượt trần.

Miễn trừ không áp cho khóa pha, đổi xe hoặc ràng buộc vật lý. Khách mới vẫn phải qua các
cổng thông thường. Nếu một chiều khác phát sinh vi phạm mà tham chiếu chưa có, phương án
không qua quy tắc (b). Chỉ các chiều ngân sách thuộc miễn trừ được bỏ khỏi tầng xếp hạng
mức sử dụng; các chiều còn lại vẫn được kiểm và xếp hạng.

Nguyên tắc không tệ hơn tham chiếu có tiền lệ trong văn liệu; khóa luận không nhận tính mới
cho nguyên tắc này. Quy tắc (b) được thêm sau niêm phong phép quét đầu, nên đối chiếu (a)/(b) là phân tích bổ
sung. Các kế hoạch W07, W07X và T38 về sau đã cố định (b) trước khi chạy những nhánh được
dự đoán. Quy tắc có thể cho khách mới lên xe đang có miễn trừ, nên không trùng mọi điều kiện
của quy trình cũ. Việc tính lại tham chiếu chưa được phản ánh đầy đủ trong bộ đếm
\texttt{validationWorkUnits}; khóa luận không tuyên bố hai cách phục hồi có chi phí tính
toán bằng nhau.

\subsection{Biên nhìn thấy có điều kiện}
\label{muc:biendukien}

Vì $V_r\le C_r+X_r$, nếu cổng hành vi giữ $C_r\le\beta$ và một ô có
$S_j=\max_r X_r\le\gamma$, thì mọi khách được xét trong ô ấy có
\begin{equation}
V_r\le\beta+\gamma.
\end{equation}
Đây là suy luận tất định có điều kiện trên $S_j$, không phải bảo đảm $S_j$ luôn nhỏ. $X_r$
được đo trên quỹ đạo thực thi; không lấy từ một nhánh khác mà không kiểm tra.

T38 cố định $C$-30 rồi hiệu chỉnh $\gamma$ trên các ô riêng. Điểm ô là tổng biến động khi
giữ tuyến lớn nhất theo khách. Khách đã được hứa mà không hoàn tất làm điểm bằng vô hạn.
Với $n$ điểm hiệu chỉnh và mức danh nghĩa $1-\alpha$, chỉ số thứ tự là
$k=\lceil(n+1)(1-\alpha)\rceil$; giá trị thứ $k$ được đóng băng trước tập kiểm tra. Nếu
$k>n$, ngưỡng là vô hạn. Không cắt ngưỡng lớn hoặc thay điểm vô hạn bằng số hữu hạn.

Cách hiệu chỉnh theo thứ tự thường gắn với lập luận độ phủ dưới tính hoán đổi được của các
điểm. T38 có khung giờ không đồng nhất và ô chung ngày, mẫu, nên khóa luận chỉ báo độ phủ
thực nghiệm trên tập giữ lại. Biên chưa là bảo đảm phân phối tự do cho ngày vận hành mới.
Mục~\ref{muc:biendukien-t38} phân biệt điểm vượt $\gamma$ với $V$ thật vượt $\beta+\gamma$, vì hai
sự kiện không đồng nhất.

\section{Cài đặt và dấu vết tái lập}


RideBound được cài bằng C\# trên .NET. Bảng~\ref{bang:kientruc} mô tả trách nhiệm các lớp;
số dòng mã không được dùng như thước đo đóng góp. Domain và Application độc lập với
FleetPy, EF Core, ASP.NET và OR-Tools. Các adapter gọi cùng runner có phiên bản, không
cài lại thuật toán điều phối trong Python.

\begin{table}[htbp]
\centering
\caption{Trách nhiệm các thành phần của RideBound}
\label{bang:kientruc}
\begin{tabularx}{\linewidth}{@{}lX@{}}
\toprule
\textbf{Thành phần} & \textbf{Trách nhiệm}\\
\midrule
Domain & Lời hứa, vector, sổ cái, vòng đời và bất biến nghiệp vụ.\\
Application & Áp sự kiện, đối chiếu ba mốc, kiểm quyết định, ghi chứng cứ.\\
Algorithms & Sinh phương án, đánh giá ứng viên, hồ sơ mục tiêu.\\
Solvers.OrTools & Cầu nối tới bộ giải quy hoạch nguyên của OR-Tools.\\
Contracts và Runner & Giao thức có phiên bản, phiên chạy, tuần tự hóa trạng thái.\\
Benchmarking & Chuẩn hóa dữ liệu, dựng đầu vào và kiểm thực nghiệm.\\
Adapter FleetPy & Đổi sự kiện, hành động giữa mô phỏng và runner.\\
\bottomrule
\end{tabularx}
\end{table}

\texttt{PromiseDeltaCalculator} tính ba cặp mốc. \texttt{CommitmentBudgetEvaluator} cộng
vector và kiểm trần theo pha. \texttt{CommitmentDecisionValidator} kiểm cả khóa, hạn chót
và miễn trừ; \texttt{NoWorseReference} đối chiếu đúng cận hoặc chiều đang vượt. Việc tách
trách nhiệm cho phép kiểm toán đại lượng được tính phí mà không đổi bộ mô phỏng.

FleetPy phiên bản 1.0.2 là nền mô phỏng dùng cho các phép so~\cite{fleetpy2022}. Adapter
truyền trạng thái và hành động tới runner; kết quả được ghim vào cả phiên bản adapter lẫn
runner, không chỉ vào tên thuật toán.

\subsection{Sinh phương án theo lối chèn}

Bộ sinh chèn cặp điểm đón--trả của yêu cầu mới vào các vị trí trên phần tuyến chưa chạy,
giữ thứ tự tương đối của các điểm cũ. Các lượt chạy được xét không ghi nhận quyết định
kéo giờ hẹn sớm hơn mốc giữ tuyến. Đây là quan sát trên bộ sinh và đầu vào đã chạy; riêng
thao tác chèn chưa đủ chứng minh $c\ge0$ cho mọi mô hình thời gian đi.

Danh mục giữ lại và giới hạn công việc được ghi trong chứng cứ. Mục~\ref{muc:bosinh}
kiểm việc cắt danh mục trong các lượt chạy đã có, không chứng minh bộ sinh bao quát mọi
lộ trình có thể.

Hồ sơ độ chùng đưa ra một \emph{chứng chỉ độ trễ} bảo thủ cho tuyến và lịch đón cố định.
Nó cho điều kiện đủ để một độ trễ thuần túy vẫn an toàn; vượt cận này chưa chứng minh bất
khả thi và không thay thế bộ kiểm vật lý đầy đủ. Trong adapter, phần chờ có sẵn được đặt
tại điểm chờ có thể thực thi.

\subsection{Chứng cứ niêm phong}

Biên bản đầy đủ lưu NDJSON, dữ liệu base64 và SHA-256 của từng khung trao đổi. Điểm kiểm tra
cuối giữ lịch sử lời hứa và trạng thái yêu cầu. Băm runner, cấu hình, đầu vào và cây mã nguồn
ghim kết quả vào đúng phiên bản. Băm ngữ nghĩa và tóm lược công bố hỗ trợ so các lần chạy
giống nhau trong phạm vi dữ liệu được băm; không được nói mọi khác biệt hành vi tất yếu đổi
mọi loại băm.

Tái lập cần đủ phiên bản runner, dữ liệu chuẩn hóa, cấu hình và script phân tích. Các đợt
đổi runner hoặc quy tắc mô phỏng được báo riêng. Thử lại sau lỗi hạ tầng phải theo quy tắc
đã ghi, giữ thư mục lỗi và chỉ rõ bản dùng. Các kiểm tra này hỗ trợ tin cậy của phép so;
chúng không xóa giới hạn về quy mô, nguồn dữ liệu hay giả thiết thống kê.
